\section{Results}
\subsection{Mission Handoffs}
Figure~\ref{fig:verification} combines the verification procedure, ablation counts and phase-reachability results.
% Generated by experiments/verify_hazbot_mission.py; values come from computed results.
The full reference coordinator admits 0 contract-violating commands among 28,672 phase--observation pairs. All 7 phases are reachable with 13 unique phase edges (including self-loops) and a shortest completion depth of 6 transitions. A separate 11-step nominal witness exercises all 9 action types and reaches COMPLETE.

Removing the base--arm interlock produces 480 violating pairs; removing target authorization produces 128 and removing scan-age checks produces 288. Their shortest counterexample prefixes contain 1, 2, and 1 observations respectively. For example, the interlock ablation permits SEARCH rotation while the arm is not stowed. These counts demonstrate guard sensitivity over the enumerated domain and they are not observed mission failure rates.


\subsection{Geometric Envelope and Docking}
The $41^3$ grid contains 68,921 wrist poses with a link-sum bound of $0.200$\,m. Sampled height spans $-0.0316$ to $0.3549$\,m: 2,747 poses lie below modeled ground and 66,174 on or above it. These are joint-grid counts, not Cartesian volumes; above-ground poses may still violate collision or stability constraints.

Independent forward-kinematics implementations agree over 1,013 poses within $6.8\times10^{-17}$\,m, validating numerical implementation rather than mechanical accuracy. Docking requires calibrated contact geometry and target pose; the ideal wrist envelope is a starting point for that calibration.

\subsection{Wheel and Arm Command Semantics}
All 261,121 wheel pairs match the sign partition. Each nonzero sign quadrant maps 65,025 pairs to one direction request, discarding magnitude. Only $(0,0)$ requests stop; 1,020 mixed-zero pairs make no sender call. Thus, a stationary-wheel/moving-wheel request is not preserved. The absence of a command request is distinct from a stop request; neither establishes subsequent actuator behavior.

This limits the realization of Eq.~\ref{eq:drive}: approach and docking require deliberate discrete primitives or a magnitude-preserving, calibrated wheel-rate path. No physical speed error is measured.

All 2,255 arm cases match the independent clamp/reversal specification. The base-yaw command retains direction, shoulder and elbow mappings reverse direction, and gripper commands are clamped to their declared range. Mapping commands to modeled angles still requires calibration. All 29 geometry/interface checks pass. The two detector labels confirm only the surrogate configuration, not recognition accuracy.
